\subsection{提高选读：动量、Adam／AdamW、Warmup 与衰减}
\phantomsection\label{app:optimizer-details}

这是第三章训练工程小节的扩展内容。\textbf{初读只需掌握 SGD：}$w_{t+1}=w_t-\eta g_t$；以下手算保留给第二遍或对应的附录 C 练习。

\subsubsection{两步小算例：动量和 Adam 为什么要记住过去？}

SGD 每次使用当前梯度。若梯度有噪声，可以保存\textbf{历史梯度的滑动平均}。
先用人为选择的两步标量梯度 $g_1=2,g_2=1$ 演示。
令动量平均系数为 $\beta_1=0.5$，初值 $m_0=0$，则
\[
m_t=0.5m_{t-1}+0.5g_t,\qquad m_1=1,\quad m_2=1.
\]
这里展示的是带归一权重的指数移动平均；其他动量实现可能使用不同的缩放约定。

Adam 还跟踪\textbf{梯度平方的移动平均}。取 $\beta_2=0.5$，
初值 $v_0=0$，同样两步可得
\[
v_t=0.5v_{t-1}+0.5g_t^2,\qquad
v_1=2,\quad v_2=1.5.
\]
由于两组平均都从零开始，Adam 还用时间步 $t$ 修正初值偏差：
\[
\widehat m_t=\frac{m_t}{1-0.5^t},\qquad
\widehat v_t=\frac{v_t}{1-0.5^t}.
\]
因此
\[
(\widehat m_1,\widehat v_1)=(2,4),\qquad
(\widehat m_2,\widehat v_2)=\left(\frac43,2\right).
\]
暂不考虑很小的数值稳定项 $\varepsilon$，
以 $-\eta\widehat m_t/\sqrt{\widehat v_t}$ 调参，
第一步是 $-\eta$，第二步约为 $-0.943\eta$。
这不是生产环境的推荐系数，只是可以手算的说明：
\textbf{Adam 使用了梯度的历史平均与历史平方平均，更新幅度不只由当前梯度决定}。

\subsubsection{学习率 Warmup 与衰减：五步怎样变化？}
\phantomsection\label{concept:lr-warmup}

开始训练时，模型权重与优化器的历史统计量尚在调整。
\textbf{Warmup} 可以让学习率先从较小值逐步增加；
随后可保持平台值，再按预定规则衰减。
它调节的是\textbf{更新步长}，并不改变正文的损失函数。

取单个参数 $w_0=1$，五步都假设梯度 $g_t=2$，
仅为了看清学习率调度对 $w_t=w_{t-1}-\eta_t g_t$ 的影响。
前三步线性升温至 $\eta_{\max}=0.06$，
第四步保持，第五步减半：
\begin{center}\small
\tagpdfsetup{table/header-rows={1}}
\begin{tabularx}{\linewidth}{@{}cXcc@{}}
\toprule
步数 $t$ & 阶段 & $\eta_t$ & 更新后 $w_t$\\\midrule
1 & Warmup & $0.02$ & $0.96$\\
2 & Warmup & $0.04$ & $0.88$\\
3 & Warmup & $0.06$ & $0.76$\\
4 & 平台 & $0.06$ & $0.64$\\
5 & 衰减 & $0.03$ & $\boxed{0.58}$\\
\bottomrule
\end{tabularx}
\end{center}
这里固定梯度只是用于\textbf{隔离学习率调度}；
实际梯度会随参数与训练批次变化。
五步手算说明的是“每步走多远”，
而不是 warmup 必定提高最终准确率。

\subsubsection{AdamW：把权重衰减与梯度更新分开}
\phantomsection\label{concept:adamw}

AdamW 继承 Adam 根据梯度历史估计\textbf{自适应更新方向}的部分，
同时把权重衰减作为单独的参数缩小步骤：
\begin{equation}
w_t=(1-\eta_t\lambda)w_{t-1}
-\eta_t\frac{\widehat m_t}{\sqrt{\widehat v_t}+\varepsilon}.
\label{eq:adamw-decoupled}
\end{equation}
$\lambda$ 是\textbf{解耦权重衰减系数}；式子中不含
偏置校正的推导，$\widehat m_t,\widehat v_t$ 已表示校正后的估计。

\begin{example}
\noindent\emph{一个参数的第一步 AdamW}\par
取 $w_0=2$、$\eta=0.1$、$\lambda=0.05$，
梯度 $g_1=0.5$。第一步偏置校正后，
$\widehat m_1=0.5,\widehat v_1=0.25$；
忽略足够小的 $\varepsilon$，自适应方向是 $1$。
于是
\[
w_1=(1-0.1\times0.05)\times2-0.1\times1
=\boxed{1.89}.
\]
\textbf{拆开看：}原权重缩小 $0.01$，再沿自适应方向更新 $0.1$。
没有衰减时，同一自适应方向得到 $1.90$。
这与直接把 $\lambda w$ 塞进 Adam 梯度的做法不同；
后者会改变动量和方差的历史统计。
\end{example}

\begin{boundarybox}
Warmup、后续衰减和 AdamW 解耦权重衰减是\textbf{不同维度的选择}：
前两者规定每一步的 $\eta_t$，后者规定优化器如何使用梯度
及独立的权重缩小项。生产训练中的批次、梯度与超参数都需要另外设定。
\end{boundarybox}

\noindent\emph{对应练习：}
\hyperref[exercise:warmup-schedule]{C1-17 的三步 Warmup}和
\hyperref[exercise:adamw-decay]{C1-18 的 AdamW 解耦衰减}，
请在看答案前分别复算。

